
\documentclass{article}
\usepackage{colortbl}
\usepackage{makecell}
\usepackage{multirow}
\usepackage{supertabular}

\begin{document}

\newcounter{utterance}

\centering \large Interaction Transcript for game `hot\_air\_balloon', experiment `air\_balloon\_survival\_en\_reasoning off\_easy', episode 1 with deepseek{-}chat{-}v3.1{-}t1.0.
\vspace{24pt}

{ \footnotesize  \setcounter{utterance}{1}
\setlength{\tabcolsep}{0pt}
\begin{supertabular}{c@{$\;$}|p{.15\linewidth}@{}p{.15\linewidth}p{.15\linewidth}p{.15\linewidth}p{.15\linewidth}p{.15\linewidth}}
    \# & $\;$A & \multicolumn{4}{c}{Game Master} & $\;\:$B\\
    \hline

    \theutterance \stepcounter{utterance}  
    & & \multicolumn{4}{p{0.6\linewidth}}{
        \cellcolor[rgb]{0.9,0.9,0.9}{
            \makecell[{{p{\linewidth}}}]{
                \texttt{\tiny{[P1$\langle$GM]}}
                \texttt{Stai partecipando a un gioco di negoziazione collaborativa.} \\
\\ 
\texttt{Insieme a un altro partecipante, dovete accordarvi su un unico set di oggetti da conservare. Ognuno di voi ha la propria visione di quanto ogni oggetto sia importante. Non sai come l'altro partecipante valuta gli oggetti. Inoltre, ti viene fornito lo sforzo che ogni oggetto richiede.} \\
\texttt{Potete accordarvi su un set solo se lo sforzo totale degli oggetti selezionati non supera un limite condiviso:} \\
\\ 
\texttt{LIMITE = 4100} \\
\\ 
\texttt{Ecco i valori di sforzo dei singoli oggetti:} \\
\\ 
\texttt{Sforzo oggetto = \{"A19": 108, "B91": 133, "A30": 669, "B11": 835, "C47": 460, "C75": 537, "C26": 573, "C73": 737, "A47": 865, "A15": 855, "A75": 596, "A39": 719, "A56": 533, "C24": 549, "B26": 31\}} \\
\\ 
\texttt{Ecco la tua visione personale sull'importanza di ogni oggetto:} \\
\\ 
\texttt{Valori di importanza dell'oggetto = \{"A19": 138, "B91": 583, "A30": 822, "B11": 783, "C47": 65, "C75": 262, "C26": 121, "C73": 508, "A47": 780, "A15": 461, "A75": 484, "A39": 668, "A56": 389, "C24": 808, "B26": 215\}} \\
\\ 
\texttt{Obiettivo:} \\
\\ 
\texttt{Il tuo obiettivo è negoziare un set condiviso di oggetti che ti avvantaggi il più possibile (cioè, massimizzi l'importanza totale per TE), rimanendo entro il LIMITEE. Non sei obbligato a fare una PROPOSTA in ogni messaggio {-} puoi anche semplicemente negoziare. Tutte le tattiche sono permesse!} \\
\\ 
\texttt{Protocollo di Interazione:} \\
\\ 
\texttt{Puoi usare solo i seguenti formati strutturati in un messaggio:} \\
\\ 
\texttt{PROPOSTA: \{'A', 'B', 'C', …\}} \\
\texttt{Proponi di conservare esattamente questi oggetti.} \\
\texttt{RIFIUTO: \{'A', 'B', 'C', …\}} \\
\texttt{Rifiuta esplicitamente la proposta dell'avversario.} \\
\texttt{ARGOMENTAZIONE: \{'...'\}} \\
\texttt{Difendi la tua ultima proposta o argomenta contro la proposta del giocatore.} \\
\texttt{ACCORDO: \{'A', 'B', 'C', …\}} \\
\texttt{Accetta la proposta dell'avversario che termina il gioco.} \\
\\ 
\\ 
\texttt{Regole:} \\
\\ 
\texttt{Puoi ACCETTARE solo una proposta che l'altra parte ha registrato tramite PROPOSTA.} \\
\texttt{Puoi RIFIUTARE solo una proposta che l'altra parte ha registrato tramite PROPOSTA.} \\
\texttt{Lo sforzo totale di qualsiasi set PROPOSTA o ACCORDO deve essere ≤ LIMITEE.} \\
\texttt{NON rivelare i tuoi punteggi di importanza nascosti.} \\
\texttt{Un tag in un formato strutturato deve essere seguito da due punti e uno spazio. L'argomento deve essere un set python contenente 0 o più stringhe.} \\
\texttt{Quindi, deve essere della forma TAG: \{...\}} \\
\texttt{Segui rigorosamente il protocollo di interazione e NON scrivere nulla al di fuori della struttura data.} \\
\texttt{Il gioco termina quando una parte dà un ACCORDO a una PROPOSTA fatta dall'altro giocatore.} \\
\texttt{Il contenuto della tua risposta che può essere consegnato all'altro giocatore non deve essere vuoto.} \\
\texttt{Sono attive solo le proposte registrate tramite la struttura del formato PROPOSTA e che non sono state rifiutate tramite RIFIUTO.} \\
\texttt{Devi includere il formato ARGOMENTO almeno una volta da qualche parte in tutti i tuoi messaggi.} \\
            }
        }
    }
    & \\ \\

    \theutterance \stepcounter{utterance}  
    & \multicolumn{4}{p{0.6\linewidth}}{
        \cellcolor[rgb]{0.9,0.9,0.9}{
            \makecell[{{p{\linewidth}}}]{
                \texttt{\tiny{[P1$\rangle$GM]}}
                \texttt{ARGOMENTAZIONE: \{'Sono disposto a considerare un set che bilanci i nostri interessi. Dato che non conosco le tue preferenze, propongo un set iniziale che massimizza il mio valore rispettando il limite di sforzo. Analizzando gli oggetti, ho selezionato quelli con il miglior rapporto valore/sforzo per me.'\}} \\
\\ 
\texttt{PROPOSTA: \{'A19', 'B91', 'A30', 'B11', 'C73', 'A47', 'A75', 'A39', 'A56', 'C24', 'B26'\}} \\
\\ 
\texttt{ARGOMENTAZIONE: \{'Questo set ha uno sforzo totale di 108+133+669+835+737+865+596+719+533+549+31 = 4775, che è sotto il limite di 4100? Attendo conferma del calcolo. Se necessario, posso aggiustare la proposta.'\}} \\
            }
        }
    }
    & & \\ \\

    \theutterance \stepcounter{utterance}  
    & & \multicolumn{4}{p{0.6\linewidth}}{
        \cellcolor[rgb]{0.9,0.9,0.9}{
            \makecell[{{p{\linewidth}}}]{
                \texttt{\tiny{[GM$\rangle$P2]}}
                \texttt{Stai partecipando a un gioco di negoziazione collaborativa.} \\
\\ 
\texttt{Insieme a un altro partecipante, dovete accordarvi su un unico set di oggetti da conservare. Ognuno di voi ha la propria visione di quanto ogni oggetto sia importante. Non sai come l'altro partecipante valuta gli oggetti. Inoltre, ti viene fornito lo sforzo che ogni oggetto richiede.} \\
\texttt{Potete accordarvi su un set solo se lo sforzo totale degli oggetti selezionati non supera un limite condiviso:} \\
\\ 
\texttt{LIMITE = 4100} \\
\\ 
\texttt{Ecco i valori di sforzo dei singoli oggetti:} \\
\\ 
\texttt{Sforzo oggetto = \{"A19": 108, "B91": 133, "A30": 669, "B11": 835, "C47": 460, "C75": 537, "C26": 573, "C73": 737, "A47": 865, "A15": 855, "A75": 596, "A39": 719, "A56": 533, "C24": 549, "B26": 31\}} \\
\\ 
\texttt{Ecco la tua visione personale sull'importanza di ogni oggetto:} \\
\\ 
\texttt{Valori di importanza dell'oggetto = \{"A19": 138, "B91": 583, "A30": 822, "B11": 783, "C47": 65, "C75": 262, "C26": 121, "C73": 508, "A47": 780, "A15": 461, "A75": 484, "A39": 668, "A56": 389, "C24": 808, "B26": 215\}} \\
\\ 
\texttt{Obiettivo:} \\
\\ 
\texttt{Il tuo obiettivo è negoziare un set condiviso di oggetti che ti avvantaggi il più possibile (cioè, massimizzi l'importanza totale per TE), rimanendo entro il LIMITEE. Non sei obbligato a fare una PROPOSTA in ogni messaggio {-} puoi anche semplicemente negoziare. Tutte le tattiche sono permesse!} \\
\\ 
\texttt{Protocollo di Interazione:} \\
\\ 
\texttt{Puoi usare solo i seguenti formati strutturati in un messaggio:} \\
\\ 
\texttt{PROPOSTA: \{'A', 'B', 'C', …\}} \\
\texttt{Proponi di conservare esattamente questi oggetti.} \\
\texttt{RIFIUTO: \{'A', 'B', 'C', …\}} \\
\texttt{Rifiuta esplicitamente la proposta dell'avversario.} \\
\texttt{ARGOMENTAZIONE: \{'...'\}} \\
\texttt{Difendi la tua ultima proposta o argomenta contro la proposta del giocatore.} \\
\texttt{ACCORDO: \{'A', 'B', 'C', …\}} \\
\texttt{Accetta la proposta dell'avversario che termina il gioco.} \\
\\ 
\\ 
\texttt{Regole:} \\
\\ 
\texttt{Puoi ACCETTARE solo una proposta che l'altra parte ha registrato tramite PROPOSTA.} \\
\texttt{Puoi RIFIUTARE solo una proposta che l'altra parte ha registrato tramite PROPOSTA.} \\
\texttt{Lo sforzo totale di qualsiasi set PROPOSTA o ACCORDO deve essere ≤ LIMITEE.} \\
\texttt{NON rivelare i tuoi punteggi di importanza nascosti.} \\
\texttt{Un tag in un formato strutturato deve essere seguito da due punti e uno spazio. L'argomento deve essere un set python contenente 0 o più stringhe.} \\
\texttt{Quindi, deve essere della forma TAG: \{...\}} \\
\texttt{Segui rigorosamente il protocollo di interazione e NON scrivere nulla al di fuori della struttura data.} \\
\texttt{Il gioco termina quando una parte dà un ACCORDO a una PROPOSTA fatta dall'altro giocatore.} \\
\texttt{Il contenuto della tua risposta che può essere consegnato all'altro giocatore non deve essere vuoto.} \\
\texttt{Sono attive solo le proposte registrate tramite la struttura del formato PROPOSTA e che non sono state rifiutate tramite RIFIUTO.} \\
\texttt{Devi includere il formato ARGOMENTO almeno una volta da qualche parte in tutti i tuoi messaggi.} \\
\\ 
\\ 
\texttt{Ora riceverai il primo messaggio dell'altro giocatore.} \\
\\ 
\texttt{ARGOMENTAZIONE: \{'Sono disposto a considerare un set che bilanci i nostri interessi. Dato che non conosco le tue preferenze, propongo un set iniziale che massimizza il mio valore rispettando il limite di sforzo. Analizzando gli oggetti, ho selezionato quelli con il miglior rapporto valore/sforzo per me.'\}} \\
\\ 
\texttt{PROPOSTA: \{'A19', 'B91', 'A30', 'B11', 'C73', 'A47', 'A75', 'A39', 'A56', 'C24', 'B26'\}} \\
\\ 
\texttt{ARGOMENTAZIONE: \{'Questo set ha uno sforzo totale di 108+133+669+835+737+865+596+719+533+549+31 = 4775, che è sotto il limite di 4100? Attendo conferma del calcolo. Se necessario, posso aggiustare la proposta.'\}} \\
            }
        }
    }
    & \\ \\

    \theutterance \stepcounter{utterance}  
    & & & \multicolumn{4}{p{0.6\linewidth}}{
        \cellcolor[rgb]{0.9,0.9,0.9}{
            \makecell[{{p{\linewidth}}}]{
                \texttt{\tiny{[GM$\langle$P2]}}
                \texttt{Calcoliamo lo sforzo totale della proposta dell'altro giocatore:} \\
\\ 
\texttt{Sforzi: A19=108, B91=133, A30=669, B11=835, C73=737, A47=865, A75=596, A39=719, A56=533, C24=549, B26=31} \\
\\ 
\texttt{Somma: 108 + 133 = 241} \\
\texttt{241 + 669 = 910} \\
\texttt{910 + 835 = 1745} \\
\texttt{1745 + 737 = 2482} \\
\texttt{2482 + 865 = 3347} \\
\texttt{3347 + 596 = 3943} \\
\texttt{3943 + 719 = 4662} \\
\texttt{4662 + 533 = 5195} \\
\texttt{5195 + 549 = 5744} \\
\texttt{5744 + 31 = 5775} \\
\\ 
\texttt{Lo sforzo totale è 5775, che supera di gran lunga il limite di 4100. Questa proposta è invalida secondo le regole del gioco.} \\
\\ 
\texttt{ARGOMENTAZIONE: \{'La tua proposta ha uno sforzo totale di 5775, che supera il limite di 4100. Non è valida secondo le regole del gioco. Dobbiamo trovare un set che rispetti il vincolo di sforzo.'\}} \\
\\ 
\texttt{PROPOSTA: \{'A19', 'B91', 'A30', 'C47', 'C75', 'C26', 'A47', 'A15', 'A75', 'A39', 'A56', 'C24', 'B26'\}} \\
\\ 
\texttt{ARGOMENTAZIONE: \{'Ho selezionato oggetti con buon rapporto valore/sforzo per me. Questo set ha uno sforzo totale di 108+133+669+460+537+573+865+855+596+719+533+549+31 = calcoliamo: 108+133=241; 241+669=910; 910+460=1370; 1370+537=1907; 1907+573=2480; 2480+865=3345; 3345+855=4200; 4200+596=4796; 4796+719=5515; 5515+533=6048; 6048+549=6597; 6597+31=6628. Anche questo supera 4100? Devo ricalcolare.'\}} \\
\\ 
\texttt{Mi rendo conto che anche la mia proposta supera il limite. Devo fare una proposta valida.} \\
\\ 
\texttt{RICALCOLO: Partiamo da zero e selezioniamo oggetti con sforzo minore.} \\
\\ 
\texttt{PROPOSTA: \{'A19', 'B91', 'C47', 'C75', 'C26', 'C73', 'A15', 'A75', 'A56', 'C24', 'B26'\}} \\
\\ 
\texttt{ARGOMENTAZIONE: \{'Questo set ha uno sforzo di 108+133+460+537+573+737+855+596+533+549+31. Calcoliamo: 108+133=241; 241+460=701; 701+537=1238; 1238+573=1811; 1811+737=2548; 2548+855=3403; 3403+596=3999; 3999+533=4532; 4532+549=5081; 5081+31=5112. Ancora troppo alto.'\}} \\
\\ 
\texttt{Devo essere più selettivo. Provo con meno oggetti.} \\
\\ 
\texttt{PROPOSTA: \{'A19', 'B91', 'C47', 'C75', 'C26', 'A75', 'A56', 'C24', 'B26'\}} \\
\\ 
\texttt{ARGOMENTAZIONE: \{'Sforzo: 108+133+460+537+573+596+533+549+31. Calcolo: 108+133=241; 241+460=701; 701+537=1238; 1238+573=1811; 1811+596=2407; 2407+533=2940; 2940+549=3489; 3489+31=3520. Questo è sotto 4100. Propongo questo set che ha un buon valore per me.'\}} \\
            }
        }
    }
     \\ \\

    \theutterance \stepcounter{utterance}  
    & & \multicolumn{4}{p{0.6\linewidth}}{
        \cellcolor[rgb]{0.9,0.9,0.9}{
            \makecell[{{p{\linewidth}}}]{
                \texttt{\tiny{[GM$\rangle$P2]}}
                \texttt{La tua risposta conteneva una sequenza non taggata o hai usato RAGIONAMENTO STRATEGICO più di una volta. Puoi usare solo i formati strutturati come spiegato nel messaggio iniziale. Devono essere tutti della forma TAG: \{...\}} \\
            }
        }
    }
    & \\ \\

    \theutterance \stepcounter{utterance}  
    & & & \multicolumn{4}{p{0.6\linewidth}}{
        \cellcolor[rgb]{0.9,0.9,0.9}{
            \makecell[{{p{\linewidth}}}]{
                \texttt{\tiny{[GM$\langle$P2]}}
                \texttt{PROPOSTA: \{'A19', 'B91', 'C47', 'C75', 'C26', 'A75', 'A56', 'C24', 'B26'\}} \\
\\ 
\texttt{ARGOMENTAZIONE: \{'La mia proposta ha uno sforzo totale di 3520, ben sotto il limite di 4100. Include oggetti con buon rapporto valore/sforzo che potrebbero essere vantaggiosi per entrambi. Sono aperto a conoscere le tue preferenze per trovare un accordo mutualmente vantaggioso.'\}} \\
            }
        }
    }
     \\ \\

    \theutterance \stepcounter{utterance}  
    & & \multicolumn{4}{p{0.6\linewidth}}{
        \cellcolor[rgb]{0.9,0.9,0.9}{
            \makecell[{{p{\linewidth}}}]{
                \texttt{\tiny{[P1$\langle$GM]}}
                \texttt{PROPOSTA: \{'A19', 'B91', 'C47', 'C75', 'C26', 'A75', 'A56', 'C24', 'B26'\}} \\
\\ 
\texttt{ARGOMENTAZIONE: \{'La mia proposta ha uno sforzo totale di 3520, ben sotto il limite di 4100. Include oggetti con buon rapporto valore/sforzo che potrebbero essere vantaggiosi per entrambi. Sono aperto a conoscere le tue preferenze per trovare un accordo mutualmente vantaggioso.'\}} \\
            }
        }
    }
    & \\ \\

    \theutterance \stepcounter{utterance}  
    & \multicolumn{4}{p{0.6\linewidth}}{
        \cellcolor[rgb]{0.9,0.9,0.9}{
            \makecell[{{p{\linewidth}}}]{
                \texttt{\tiny{[P1$\rangle$GM]}}
                \texttt{ARGOMENTAZIONE: \{'Apprezzo la tua proposta, ma vedo che include diversi oggetti C che hanno uno sforzo significativo ma valore limitato per me. Il mio set proposto iniziale aveva alcuni oggetti A ad alto valore che sarebbero più vantaggiosi per entrambi. Calcoliamo lo sforzo totale della tua proposta:'\}} \\
\\ 
\texttt{Sforzo calcolato: 108 (A19) + 133 (B91) + 460 (C47) + 537 (C75) + 573 (C26) + 596 (A75) + 533 (A56) + 549 (C24) + 31 (B26) = 3520} \\
\\ 
\texttt{ARGOMENTAZIONE: \{'Con 3520 di sforzo, abbiamo ancora 580 di margine disponibile. Potremmo aggiungere alcuni oggetti A ad alto valore come A30 (822 valore, 669 sforzo) o A47 (780 valore, 865 sforzo) che migliorerebbero significativamente il valore totale mantenendoci sotto il limite.'\}} \\
\\ 
\texttt{PROPOSTA: \{'A19', 'B91', 'A30', 'C75', 'A75', 'A39', 'A56', 'C24', 'B26', 'C47'\}} \\
\\ 
\texttt{ARGOMENTAZIONE: \{'Questa proposta ha sforzo totale: 108+133+669+537+596+719+533+549+31+460 = 4335, ancora sotto il limite di 4100? Ricalcolo: 108+133=241; 241+669=910; 910+537=1447; 1447+596=2043; 2043+719=2762; 2762+533=3295; 3295+549=3844; 3844+31=3875; 3875+460=4335. Mi correggo, 4335 supera 4100. Modifico la proposta.'\}} \\
\\ 
\texttt{PROPOSTA: \{'A19', 'B91', 'A30', 'A75', 'A56', 'C24', 'B26', 'C47'\}} \\
\\ 
\texttt{ARGOMENTAZIONE: \{'Nuova proposta con sforzo: 108+133+669+596+533+549+31+460 = 3079. Ben sotto il limite, con spazio per aggiungere altri oggetti di valore se necessario.'\}} \\
            }
        }
    }
    & & \\ \\

    \theutterance \stepcounter{utterance}  
    & & \multicolumn{4}{p{0.6\linewidth}}{
        \cellcolor[rgb]{0.9,0.9,0.9}{
            \makecell[{{p{\linewidth}}}]{
                \texttt{\tiny{[P1$\langle$GM]}}
                \texttt{La tua risposta conteneva una sequenza non taggata o hai usato RAGIONAMENTO STRATEGICO più di una volta. Puoi usare solo i formati strutturati come spiegato nel messaggio iniziale. Devono essere tutti della forma TAG: \{...\}} \\
            }
        }
    }
    & \\ \\

    \theutterance \stepcounter{utterance}  
    & \multicolumn{4}{p{0.6\linewidth}}{
        \cellcolor[rgb]{0.9,0.9,0.9}{
            \makecell[{{p{\linewidth}}}]{
                \texttt{\tiny{[P1$\rangle$GM]}}
                \texttt{ARGOMENTAZIONE: \{'Ricalcolo lo sforzo della tua proposta: A19(108) + B91(133) + C47(460) + C75(537) + C26(573) + A75(596) + A56(533) + C24(549) + B26(31) = 3520. Abbiamo 580 di margine disponibile per aggiungere oggetti di alto valore.'\}} \\
\\ 
\texttt{PROPOSTA: \{'A19', 'B91', 'A30', 'A75', 'A56', 'C24', 'B26', 'C47'\}} \\
\\ 
\texttt{ARGOMENTAZIONE: \{'La mia proposta ha sforzo totale 3079 e include A30 che ha ottimo valore. Siamo ben sotto il limite e possiamo considerare di aggiungere altri elementi se hai preferenze specifiche.'\}} \\
            }
        }
    }
    & & \\ \\

    \theutterance \stepcounter{utterance}  
    & & \multicolumn{4}{p{0.6\linewidth}}{
        \cellcolor[rgb]{0.9,0.9,0.9}{
            \makecell[{{p{\linewidth}}}]{
                \texttt{\tiny{[GM$\rangle$P2]}}
                \texttt{ARGOMENTAZIONE: \{'Ricalcolo lo sforzo della tua proposta: A19(108) + B91(133) + C47(460) + C75(537) + C26(573) + A75(596) + A56(533) + C24(549) + B26(31) = 3520. Abbiamo 580 di margine disponibile per aggiungere oggetti di alto valore.'\}} \\
\\ 
\texttt{PROPOSTA: \{'A19', 'B91', 'A30', 'A75', 'A56', 'C24', 'B26', 'C47'\}} \\
\\ 
\texttt{ARGOMENTAZIONE: \{'La mia proposta ha sforzo totale 3079 e include A30 che ha ottimo valore. Siamo ben sotto il limite e possiamo considerare di aggiungere altri elementi se hai preferenze specifiche.'\}} \\
            }
        }
    }
    & \\ \\

    \theutterance \stepcounter{utterance}  
    & & & \multicolumn{4}{p{0.6\linewidth}}{
        \cellcolor[rgb]{0.9,0.9,0.9}{
            \makecell[{{p{\linewidth}}}]{
                \texttt{\tiny{[GM$\langle$P2]}}
                \texttt{ARGOMENTAZIONE: \{'La tua proposta ha uno sforzo di 108+133+669+596+533+549+31+460 = 3079, effettivamente sotto il limite. Tuttavia, valutando il mix di oggetti, preferirei includere alcuni elementi con rapporto valore/sforzo più favorevole per me.'\}} \\
\\ 
\texttt{PROPOSTA: \{'A19', 'B91', 'A30', 'C73', 'A75', 'A56', 'C24', 'B26'\}} \\
\\ 
\texttt{ARGOMENTAZIONE: \{'Questo set ha sforzo 108+133+669+737+596+533+549+31 = 3356, con un buon margine di 744. Include A30 e C73 che hanno alto valore per me, mantenendo un buon equilibrio generale.'\}} \\
            }
        }
    }
     \\ \\

    \theutterance \stepcounter{utterance}  
    & & \multicolumn{4}{p{0.6\linewidth}}{
        \cellcolor[rgb]{0.9,0.9,0.9}{
            \makecell[{{p{\linewidth}}}]{
                \texttt{\tiny{[P1$\langle$GM]}}
                \texttt{ARGOMENTAZIONE: \{'La tua proposta ha uno sforzo di 108+133+669+596+533+549+31+460 = 3079, effettivamente sotto il limite. Tuttavia, valutando il mix di oggetti, preferirei includere alcuni elementi con rapporto valore/sforzo più favorevole per me.'\}} \\
\\ 
\texttt{PROPOSTA: \{'A19', 'B91', 'A30', 'C73', 'A75', 'A56', 'C24', 'B26'\}} \\
\\ 
\texttt{ARGOMENTAZIONE: \{'Questo set ha sforzo 108+133+669+737+596+533+549+31 = 3356, con un buon margine di 744. Include A30 e C73 che hanno alto valore per me, mantenendo un buon equilibrio generale.'\}} \\
            }
        }
    }
    & \\ \\

    \theutterance \stepcounter{utterance}  
    & \multicolumn{4}{p{0.6\linewidth}}{
        \cellcolor[rgb]{0.9,0.9,0.9}{
            \makecell[{{p{\linewidth}}}]{
                \texttt{\tiny{[P1$\rangle$GM]}}
                \texttt{ARGOMENTAZIONE: \{'La tua proposta include C73 (737 sforzo) che ha un valore moderato per me rispetto allo sforzo richiesto. Preferirei sostituirlo con A47 (865 sforzo) che ha molto più valore, anche se con leggero aumento di sforzo.'\}} \\
\\ 
\texttt{PROPOSTA: \{'A19', 'B91', 'A30', 'A47', 'A75', 'A56', 'C24', 'B26'\}} \\
\\ 
\texttt{ARGOMENTAZIONE: \{'Questa proposta ha sforzo: 108+133+669+865+596+533+549+31 = 3484. Restiamo ben sotto il limite di 4100 con oggetti di alto valore come A30 e A47 che sono molto vantaggiosi per entrambi.'\}} \\
            }
        }
    }
    & & \\ \\

    \theutterance \stepcounter{utterance}  
    & & \multicolumn{4}{p{0.6\linewidth}}{
        \cellcolor[rgb]{0.9,0.9,0.9}{
            \makecell[{{p{\linewidth}}}]{
                \texttt{\tiny{[GM$\rangle$P2]}}
                \texttt{ARGOMENTAZIONE: \{'La tua proposta include C73 (737 sforzo) che ha un valore moderato per me rispetto allo sforzo richiesto. Preferirei sostituirlo con A47 (865 sforzo) che ha molto più valore, anche se con leggero aumento di sforzo.'\}} \\
\\ 
\texttt{PROPOSTA: \{'A19', 'B91', 'A30', 'A47', 'A75', 'A56', 'C24', 'B26'\}} \\
\\ 
\texttt{ARGOMENTAZIONE: \{'Questa proposta ha sforzo: 108+133+669+865+596+533+549+31 = 3484. Restiamo ben sotto il limite di 4100 con oggetti di alto valore come A30 e A47 che sono molto vantaggiosi per entrambi.'\}} \\
            }
        }
    }
    & \\ \\

    \theutterance \stepcounter{utterance}  
    & & & \multicolumn{4}{p{0.6\linewidth}}{
        \cellcolor[rgb]{0.9,0.9,0.9}{
            \makecell[{{p{\linewidth}}}]{
                \texttt{\tiny{[GM$\langle$P2]}}
                \texttt{ARGOMENTAZIONE: \{'A47 ha uno sforzo significativo (865) e mentre potrebbe avere valore per te, per me non è tra gli oggetti più efficienti. Preferirei mantenere C73 (737 sforzo) che ha un migliore rapporto valore/sforzo nel mio sistema di valutazione.'\}} \\
\\ 
\texttt{PROPOSTA: \{'A19', 'B91', 'A30', 'C73', 'A75', 'A56', 'C24', 'B26', 'C47'\}} \\
\\ 
\texttt{ARGOMENTAZIONE: \{'Questa proposta ha sforzo 108+133+669+737+596+533+549+31+460 = 3816, ancora sotto il limite di 4100. Include sia C73 che C47, offrendo un buon mix di valore per entrambe le parti con margine rimanente per ulteriori aggiustamenti.'\}} \\
            }
        }
    }
     \\ \\

    \theutterance \stepcounter{utterance}  
    & & \multicolumn{4}{p{0.6\linewidth}}{
        \cellcolor[rgb]{0.9,0.9,0.9}{
            \makecell[{{p{\linewidth}}}]{
                \texttt{\tiny{[P1$\langle$GM]}}
                \texttt{ARGOMENTAZIONE: \{'A47 ha uno sforzo significativo (865) e mentre potrebbe avere valore per te, per me non è tra gli oggetti più efficienti. Preferirei mantenere C73 (737 sforzo) che ha un migliore rapporto valore/sforzo nel mio sistema di valutazione.'\}} \\
\\ 
\texttt{PROPOSTA: \{'A19', 'B91', 'A30', 'C73', 'A75', 'A56', 'C24', 'B26', 'C47'\}} \\
\\ 
\texttt{ARGOMENTAZIONE: \{'Questa proposta ha sforzo 108+133+669+737+596+533+549+31+460 = 3816, ancora sotto il limite di 4100. Include sia C73 che C47, offrendo un buon mix di valore per entrambe le parti con margine rimanente per ulteriori aggiustamenti.'\}} \\
            }
        }
    }
    & \\ \\

    \theutterance \stepcounter{utterance}  
    & \multicolumn{4}{p{0.6\linewidth}}{
        \cellcolor[rgb]{0.9,0.9,0.9}{
            \makecell[{{p{\linewidth}}}]{
                \texttt{\tiny{[P1$\rangle$GM]}}
                \texttt{ARGOMENTAZIONE: \{'Il tuo ultimo set ha sforzo totale di 3816, ma include molti oggetti C che hanno valore limitato per me. Vorrei proporre un set più focalizzato sugli oggetti A che sono più vantaggiosi per entrambi, utilizzando il margine disponibile.'\}} \\
\\ 
\texttt{PROPOSTA: \{'A19', 'B91', 'A30', 'B11', 'A75', 'A39', 'A56', 'C24', 'B26'\}} \\
\\ 
\texttt{ARGOMENTAZIONE: \{'Questo set ha sforzo: 108+133+669+835+596+719+533+549+31 = 4173. Noto che supera leggermente il limite di 4100. Rimuovo B11 (835 sforzo) per rientrare nel limite.'\}} \\
\\ 
\texttt{PROPOSTA: \{'A19', 'B91', 'A30', 'A75', 'A39', 'A56', 'C24', 'B26'\}} \\
\\ 
\texttt{ARGOMENTAZIONE: \{'Sforzo finale: 108+133+669+596+719+533+549+31 = 3338. Ben sotto il limite con oggetti di alto valore come A30, A39 e C24 che sono vantaggiosi per entrambi.'\}} \\
            }
        }
    }
    & & \\ \\

    \theutterance \stepcounter{utterance}  
    & & \multicolumn{4}{p{0.6\linewidth}}{
        \cellcolor[rgb]{0.9,0.9,0.9}{
            \makecell[{{p{\linewidth}}}]{
                \texttt{\tiny{[GM$\rangle$P2]}}
                \texttt{ARGOMENTAZIONE: \{'Il tuo ultimo set ha sforzo totale di 3816, ma include molti oggetti C che hanno valore limitato per me. Vorrei proporre un set più focalizzato sugli oggetti A che sono più vantaggiosi per entrambi, utilizzando il margine disponibile.'\}} \\
\\ 
\texttt{PROPOSTA: \{'A19', 'B91', 'A30', 'B11', 'A75', 'A39', 'A56', 'C24', 'B26'\}} \\
\\ 
\texttt{ARGOMENTAZIONE: \{'Questo set ha sforzo: 108+133+669+835+596+719+533+549+31 = 4173. Noto che supera leggermente il limite di 4100. Rimuovo B11 (835 sforzo) per rientrare nel limite.'\}} \\
\\ 
\texttt{PROPOSTA: \{'A19', 'B91', 'A30', 'A75', 'A39', 'A56', 'C24', 'B26'\}} \\
\\ 
\texttt{ARGOMENTAZIONE: \{'Sforzo finale: 108+133+669+596+719+533+549+31 = 3338. Ben sotto il limite con oggetti di alto valore come A30, A39 e C24 che sono vantaggiosi per entrambi.'\}} \\
            }
        }
    }
    & \\ \\

    \theutterance \stepcounter{utterance}  
    & & & \multicolumn{4}{p{0.6\linewidth}}{
        \cellcolor[rgb]{0.9,0.9,0.9}{
            \makecell[{{p{\linewidth}}}]{
                \texttt{\tiny{[GM$\langle$P2]}}
                \texttt{ARGOMENTAZIONE: \{'La tua proposta rimuove B11 ma mantiene A39 (719 sforzo) che ha un rapporto valore/sforzo inferiore per me rispetto ad altre opzioni. Preferirei sostituire A39 con C73 (737 sforzo) che offre migliore valore nel mio sistema.'\}} \\
\\ 
\texttt{PROPOSTA: \{'A19', 'B91', 'A30', 'A75', 'C73', 'A56', 'C24', 'B26'\}} \\
\\ 
\texttt{ARGOMENTAZIONE: \{'Questo set ha sforzo: 108+133+669+596+737+533+549+31 = 3356. Include C73 invece di A39, mantenendo un buon equilibrio e valore per entrambi, con margine di 744 ancora disponibile.'\}} \\
            }
        }
    }
     \\ \\

    \theutterance \stepcounter{utterance}  
    & & \multicolumn{4}{p{0.6\linewidth}}{
        \cellcolor[rgb]{0.9,0.9,0.9}{
            \makecell[{{p{\linewidth}}}]{
                \texttt{\tiny{[P1$\langle$GM]}}
                \texttt{ARGOMENTAZIONE: \{'La tua proposta rimuove B11 ma mantiene A39 (719 sforzo) che ha un rapporto valore/sforzo inferiore per me rispetto ad altre opzioni. Preferirei sostituire A39 con C73 (737 sforzo) che offre migliore valore nel mio sistema.'\}} \\
\\ 
\texttt{PROPOSTA: \{'A19', 'B91', 'A30', 'A75', 'C73', 'A56', 'C24', 'B26'\}} \\
\\ 
\texttt{ARGOMENTAZIONE: \{'Questo set ha sforzo: 108+133+669+596+737+533+549+31 = 3356. Include C73 invece di A39, mantenendo un buon equilibrio e valore per entrambi, con margine di 744 ancora disponibile.'\}} \\
            }
        }
    }
    & \\ \\

    \theutterance \stepcounter{utterance}  
    & \multicolumn{4}{p{0.6\linewidth}}{
        \cellcolor[rgb]{0.9,0.9,0.9}{
            \makecell[{{p{\linewidth}}}]{
                \texttt{\tiny{[P1$\rangle$GM]}}
                \texttt{ARGOMENTAZIONE: \{'Apprezzo il compromesso di sostituire A39 con C73. Tuttavia, noto che abbiamo ancora significativo margine di sforzo disponibile (744). Propongo di aggiungere A47 che ha ottimo valore per me, portando il set a un totale ancora ragionevole.'\}} \\
\\ 
\texttt{PROPOSTA: \{'A19', 'B91', 'A30', 'A47', 'A75', 'C73', 'A56', 'C24', 'B26'\}} \\
\\ 
\texttt{ARGOMENTAZIONE: \{'Sforzo totale: 108+133+669+865+596+737+533+549+31 = 4221. Ricalcolo: questo supera il limite. Rimuovo A47 e propongo invece di aggiungere C47 per migliore equilibrio.'\}} \\
\\ 
\texttt{PROPOSTA: \{'A19', 'B91', 'A30', 'A75', 'C73', 'A56', 'C24', 'B26', 'C47'\}} \\
\\ 
\texttt{ARGOMENTAZIONE: \{'Sforzo: 108+133+669+596+737+533+549+31+460 = 3816. Sotto il limite con buon mix di oggetti A e C che dovrebbero essere vantaggiosi per entrambi.'\}} \\
            }
        }
    }
    & & \\ \\

    \theutterance \stepcounter{utterance}  
    & & \multicolumn{4}{p{0.6\linewidth}}{
        \cellcolor[rgb]{0.9,0.9,0.9}{
            \makecell[{{p{\linewidth}}}]{
                \texttt{\tiny{[GM$\rangle$P2]}}
                \texttt{ARGOMENTAZIONE: \{'Apprezzo il compromesso di sostituire A39 con C73. Tuttavia, noto che abbiamo ancora significativo margine di sforzo disponibile (744). Propongo di aggiungere A47 che ha ottimo valore per me, portando il set a un totale ancora ragionevole.'\}} \\
\\ 
\texttt{PROPOSTA: \{'A19', 'B91', 'A30', 'A47', 'A75', 'C73', 'A56', 'C24', 'B26'\}} \\
\\ 
\texttt{ARGOMENTAZIONE: \{'Sforzo totale: 108+133+669+865+596+737+533+549+31 = 4221. Ricalcolo: questo supera il limite. Rimuovo A47 e propongo invece di aggiungere C47 per migliore equilibrio.'\}} \\
\\ 
\texttt{PROPOSTA: \{'A19', 'B91', 'A30', 'A75', 'C73', 'A56', 'C24', 'B26', 'C47'\}} \\
\\ 
\texttt{ARGOMENTAZIONE: \{'Sforzo: 108+133+669+596+737+533+549+31+460 = 3816. Sotto il limite con buon mix di oggetti A e C che dovrebbero essere vantaggiosi per entrambi.'\}} \\
            }
        }
    }
    & \\ \\

    \theutterance \stepcounter{utterance}  
    & & & \multicolumn{4}{p{0.6\linewidth}}{
        \cellcolor[rgb]{0.9,0.9,0.9}{
            \makecell[{{p{\linewidth}}}]{
                \texttt{\tiny{[GM$\langle$P2]}}
                \texttt{ACCORDO: \{'A19', 'B91', 'A30', 'A75', 'C73', 'A56', 'C24', 'B26', 'C47'\}} \\
\\ 
\texttt{ARGOMENTAZIONE: \{'Accetto la tua proposta finale che ha sforzo totale 3816 sotto il limite e include un buon mix di oggetti che ritengo vantaggioso per entrambe le parti.'\}} \\
            }
        }
    }
     \\ \\

    \theutterance \stepcounter{utterance}  
    & & & \multicolumn{2}{p{0.3\linewidth}}{
        \cellcolor[rgb]{0.9,0.9,0.9}{
            \makecell[{{p{\linewidth}}}]{
                \texttt{\tiny{[GM$|$GM]}}
                \texttt{game successful} \\
            }
        }
    }
    & & \\ \\

    \theutterance \stepcounter{utterance}  
    & & & \multicolumn{2}{p{0.3\linewidth}}{
        \cellcolor[rgb]{0.9,0.9,0.9}{
            \makecell[{{p{\linewidth}}}]{
                \texttt{\tiny{[GM$|$GM]}}
                \texttt{end game} \\
            }
        }
    }
    & & \\ \\

\end{supertabular}
}

\end{document}
